% Self-contained single-file LaTeX source for Overleaf upload.
%
% USAGE: Create a new Overleaf project and:
%   1. Paste this entire file into main.tex.
%   2. Upload the figures from results/ as a flat list in the project root:
%        runtime_by_family.pdf  (or .png)
%        cwz_scaling.pdf
%        nsp_existence.pdf
%        adversarial.pdf
%      Without these, the four \includegraphics in chapter 5 will fail; you
%      can comment those figure blocks out if you only want the prose.
%   3. Set the compiler to "pdfLaTeX" and BibTeX (Menu -> Settings).
%
% The bibliography is embedded via filecontents* below, so no separate
% refs.bib upload is needed.
%
% Source repo: see thesis/ in the project tree.

\begin{filecontents*}{refs.bib}
@misc{ChenWeinZhang2025,
  author        = {Chen, Kai and Wein, Nicole and Zhang, Yinzhan},
  title         = {A Polynomial-Time Algorithm for the Next-to-Shortest Path Problem
                   on Positively Weighted Directed Graphs},
  year          = {2025},
  eprint        = {2511.04345},
  archiveprefix = {arXiv},
  primaryclass  = {cs.DS}
}
@article{Yen1971,
  author  = {Yen, Jin Y.},
  title   = {Finding the {K} Shortest Loopless Paths in a Network},
  journal = {Management Science},
  volume  = {17},
  number  = {11},
  pages   = {712--716},
  year    = {1971}
}
@article{Dijkstra1959,
  author  = {Dijkstra, Edsger W.},
  title   = {A note on two problems in connexion with graphs},
  journal = {Numerische Mathematik},
  volume  = {1},
  pages   = {269--271},
  year    = {1959}
}
@article{FredmanTarjan1987,
  author  = {Fredman, Michael L. and Tarjan, Robert Endre},
  title   = {Fibonacci heaps and their uses in improved network optimization algorithms},
  journal = {Journal of the ACM},
  volume  = {34},
  number  = {3},
  pages   = {596--615},
  year    = {1987}
}
@misc{LallaSrinivasan1992,
  author = {Lalla, E. and Srinivasan, S.},
  title  = {On the NP-hardness of the next-to-shortest path problem},
  year   = {1992},
  note   = {Citation needs verification before final submission.}
}
@article{KrasikovNoble2004,
  author  = {Krasikov, Ilia and Noble, Steven D.},
  title   = {Finding next-to-shortest paths in a graph},
  journal = {Information Processing Letters},
  volume  = {92},
  number  = {3},
  pages   = {117--119},
  year    = {2004}
}
@article{Wu2010,
  author  = {Wu, Bang Ye},
  title   = {A polynomial time algorithm for the next-to-shortest path problem
             on undirected positively weighted graphs},
  journal = {Information Processing Letters},
  year    = {2010},
  note    = {Confirm exact venue/year before final submission.}
}
@article{Eppstein1998,
  author  = {Eppstein, David},
  title   = {Finding the {$k$} shortest paths},
  journal = {SIAM Journal on Computing},
  volume  = {28},
  number  = {2},
  pages   = {652--673},
  year    = {1998}
}
@article{FortuneHopcroftWyllie1980,
  author  = {Fortune, Steven and Hopcroft, John E. and Wyllie, James},
  title   = {The directed subgraph homeomorphism problem},
  journal = {Theoretical Computer Science},
  volume  = {10},
  number  = {2},
  pages   = {111--121},
  year    = {1980}
}
@misc{Tholey2012,
  author = {Tholey, Torsten},
  title  = {Linear time algorithms for two disjoint paths problems on
            directed acyclic graphs},
  year   = {2012},
  note   = {Confirm exact venue (journal vs conference) before final submission.}
}
\end{filecontents*}

\documentclass[12pt,a4paper]{report}

\usepackage[T1]{fontenc}
\usepackage[utf8]{inputenc}
\usepackage[english]{babel}
\usepackage{lmodern}
\usepackage{microtype}
\usepackage{amsmath, amssymb, amsthm}
\usepackage{algorithm, algpseudocode}
\usepackage{graphicx}
\usepackage{booktabs}
\usepackage{hyperref}
\usepackage{cleveref}
\usepackage[a4paper, margin=2.5cm]{geometry}

\newtheorem{theorem}{Theorem}[chapter]
\newtheorem{lemma}[theorem]{Lemma}
\newtheorem{definition}[theorem]{Definition}

\title{Implementation and Performance Analysis of a Polynomial-Time Algorithm
       for the Next-to-Shortest Path Problem in Directed Weighted Graphs}
\author{Ihar Maroz \\ Supervisor: Lech Duraj \\ Jagiellonian University}
\date{2026}

\begin{document}

\maketitle

\begin{abstract}
TODO: write final abstract once experiments are finalised. Briefly: we
implement the Chen--Wein--Zhang (2025) polynomial-time algorithm for the
next-to-shortest path problem on positively weighted directed graphs --- the
first polynomial-time algorithm for this problem after roughly thirty years
of being open --- and empirically evaluate its practical performance
against a brute-force oracle and Yen's $k$-shortest-simple-paths heuristic.
We confirm Yen dominates on benign random inputs but demonstrate, via a
diamond-chain adversarial family with $2^{k}$ shortest paths, a regime where
CWZ is faster than Yen by tens of thousands of times.
\end{abstract}

\tableofcontents

% ============================================================================
% Chapter 1: Introduction
% ============================================================================
\chapter{Introduction}
\label{ch:intro}

\section{The next-to-shortest path problem}

Given a directed graph $G=(V,E,w)$ with strictly positive edge weights
$w:E\to\mathbb{R}_{>0}$ and two distinguished vertices $s,t\in V$, the
\emph{next-to-shortest path} (NSP) problem asks for a simple $s\to t$ path
whose weight is the smallest among those $s\to t$ paths that are \emph{not}
shortest paths. Equivalently, write
\[
  d_G(s,t)\;=\;\min\bigl\{w(P) : P \text{ is an } s\to t \text{ path}\bigr\}.
\]
An $s\to t$ path~$P$ is a \emph{not-shortest path} if $w(P) > d_G(s,t)$, and an
NSP is one such path of minimum weight, if one exists.

Despite its surface similarity to the shortest path problem, the NSP problem
behaves very differently. The shortest path problem in graphs with positive
weights is solved in near-linear time by Dijkstra's algorithm, but the NSP
problem becomes \emph{NP-complete} as soon as zero-weight edges are
admitted in a directed graph~\cite{LallaSrinivasan1992}, and the
\emph{strictly} positively weighted directed case had remained open for nearly
thirty years. Polynomial-time algorithms were known only for undirected
graphs~\cite{KrasikovNoble2004} and for planar directed
graphs~\cite{Wu2010}; the general directed positively-weighted case was
resolved only recently by Chen, Wein, and Zhang~\cite{ChenWeinZhang2025},
whose algorithm runs in $O\bigl(|V|^{4}|E|^{3}\log|V|\bigr)$ time.

\section{Motivation}

NSP appears as a primitive in many higher-level routing and resilience
problems. Whenever a system needs a backup path under the constraint that
the backup must remain ``simple'' (no vertex repetitions), the NSP gives the
cheapest single-failure alternative. Examples include alternative-route
suggestions in road networks, fast-rerouting in IP networks, and
$k$-shortest path enumeration as in Yen's algorithm where the second-shortest
\emph{simple} path matters when the shortest is unavailable. Yen's
algorithm~\cite{Yen1971} produces simple paths in non-decreasing weight order
but its worst-case runtime is exponential in the number of distinct shortest
paths; the Chen--Wein--Zhang (CWZ) algorithm, in contrast, runs in fixed
polynomial time regardless of how many shortest paths the graph contains.

\section{Contribution}

The CWZ paper is purely theoretical and contains no implementation or
experimental evaluation. The contributions of this thesis are:
\begin{enumerate}
  \item A C++20 implementation of the CWZ algorithm and its supporting
    reductions to $(s,t)$-straight and $(s,t)$-layered digraphs.
  \item Reference implementations of two correctness oracles --- a brute-force
    enumerator and Yen's $k$-shortest-simple-paths algorithm used as an NSP
    heuristic --- and a differential-fuzzing harness comparing them across
    random graphs from five generator families.
  \item An empirical study of the CWZ algorithm's runtime on graphs from
    several families (Erd\H{o}s--R\'enyi, random DAG, layered, grid,
    scale-free), comparing it against Yen-NSP as the practical baseline,
    and a deliberate adversarial benchmark on diamond-chain graphs that
    isolates CWZ's worst-case advantage. This is, to our knowledge, the
    first empirical evaluation of the CWZ algorithm.
\end{enumerate}

\section{Outline}

Chapter~\ref{ch:litreview} surveys prior work on NSP and related shortest-path
problems. Chapter~\ref{ch:algorithm} describes the CWZ algorithm, including
its two reductions, its key Lemma~5.3, and the role of $2$-vertex-disjoint
paths in directed acyclic graphs. Chapter~\ref{ch:implementation} discusses
implementation choices and a deliberate simplification we adopted in the
core layered phase. Chapter~\ref{ch:experiments} reports empirical results
across graph families and on the adversarial benchmark.
Chapter~\ref{ch:conclusion} concludes.

% ============================================================================
% Chapter 2: Literature Review
% ============================================================================
\chapter{Literature Review}
\label{ch:litreview}

\section{Shortest paths}

Single-source shortest paths in directed graphs with non-negative weights are
classically solved by Dijkstra's algorithm in
$O((|V|+|E|)\log|V|)$ time using a binary heap, or $O(|V|\log|V|+|E|)$ with a
Fibonacci heap~\cite{Dijkstra1959, FredmanTarjan1987}. We rely on the
former; for our typical graph sizes the asymptotic improvement of Fibonacci
heaps is not worth the constant factor.

\section{$k$-shortest simple paths}

Yen's algorithm~\cite{Yen1971} enumerates the $k$ shortest simple $s\to t$
paths in non-decreasing order of weight using $k$ Dijkstra-style searches per
``spur'' along the previously-found paths. Its runtime is
$O(kn(m+n\log n))$, which is polynomial in $k$ but Yen's algorithm is used
here in a slightly non-standard way: we invoke it as an NSP solver by
enumerating paths until one is found whose weight strictly exceeds
$d_G(s,t)$. If the graph has $r$ distinct simple shortest paths from $s$ to
$t$, this requires examining at least $r$ paths before finding one with
strictly greater cost, and $r$ can be exponential in $|V|$. In favorable
inputs Yen is extremely fast, but on adversarial graphs with many alternate
shortest paths it degenerates.

Eppstein~\cite{Eppstein1998} gives an $O((m + n\log n) + k)$ algorithm for
$k$-shortest paths that need not be simple. NSP requires simplicity and so
Eppstein's algorithm does not directly apply.

\section{Next-to-shortest path: positive results}

Lalla and Srinivasan~\cite{LallaSrinivasan1992} introduced the NSP problem
and showed it is NP-complete in directed graphs with non-negative weights (in
particular, when zero weights are permitted). The undirected case admits a
polynomial-time algorithm: Krasikov and Noble~\cite{KrasikovNoble2004}
gave an $O(n^3 m)$-time algorithm; Wu's later
work~\cite{Wu2010} brings the bound down to $O(nm)$ on undirected graphs.
For \emph{planar} directed graphs with positive weights, Wu also gives a
polynomial-time algorithm.

For \emph{general} directed graphs with strictly positive weights, the
problem remained open for nearly three decades. Chen, Wein, and Zhang
\cite{ChenWeinZhang2025} resolved it in November~2025 with an
$O\bigl(|V|^{4}|E|^{3}\log|V|\bigr)$-time algorithm. The algorithm proceeds
by two reductions --- general digraph to $(s,t)$-straight to $(s,t)$-layered
--- followed by a structural enumeration over $6$-tuples of vertices and
edges that, together with a $2$-vertex-disjoint-paths subroutine in a DAG,
suffice to find any NSP if one exists.

\section{Related primitives}

The CWZ algorithm depends critically on a subroutine for finding two
vertex-disjoint paths between prescribed source-sink pairs in a directed
acyclic graph. Fortune, Hopcroft, and Wyllie~\cite{FortuneHopcroftWyllie1980}
gave the first polynomial-time algorithm for the $k$-vertex-disjoint-paths
problem in DAGs for fixed $k$. Tholey~\cite{Tholey2012} subsequently gave a
linear-time algorithm for the special case $k=2$. The CWZ paper cites Tholey
for the $O(|V|+|E|)$ bound on the inner subroutine. In the present work we
substitute a simpler max-flow-based 2-VDP that costs an extra factor of
$O(|V|+|E|)$ per invocation; this preserves correctness while raising the
overall complexity by a small polynomial factor relative to the paper's
tight bound.

% ============================================================================
% Chapter 3: Algorithm
% ============================================================================
\chapter{The Chen--Wein--Zhang Algorithm}
\label{ch:algorithm}

This chapter describes the algorithm of Chen, Wein, and
Zhang~\cite{ChenWeinZhang2025} at the level of detail required to follow the
implementation in Chapter~\ref{ch:implementation}. We follow the paper's
notation. Let $G=(V,E,w)$ be a directed graph with $w:E\to\mathbb{R}_{>0}$,
and let $s,t\in V$ be two distinguished vertices.

\section{Preliminaries}

For any vertex $v$, write $d_G(s,v)$ for the shortest distance from $s$ to
$v$ in~$G$.

\begin{definition}[Forward-edge, back-edge]
An edge $(u,v)\in E$ is a \emph{forward-edge} if
$d_G(s,u)+w(u,v)=d_G(s,v)$, and a \emph{back-edge} if
$d_G(s,u)+w(u,v)>d_G(s,v)$.
\end{definition}

Equivalently, forward edges are exactly the edges lying on some shortest
$s\to v$ path. The \emph{shortest-path DAG} of $(G,s)$ is the subgraph
consisting of all forward edges.

\begin{definition}[(s,t)-straight, (s,t)-layered]
$G$ is \emph{$(s,t)$-straight} if every vertex $v$ satisfies
$d_G(s,v)+d_G(v,t)=d_G(s,t)$, i.e.~lies on some shortest $s\to t$ path. $G$
is \emph{$(s,t)$-layered} if it is $(s,t)$-straight and additionally every
edge $(u,v)\in E$ has $d_G(s,u)\neq d_G(s,v)$ and every forward edge satisfies
$d_G(s,v) = d_G(s,u)+w(u,v)$ where the gap equals the consecutive
distance-layer spacing.
\end{definition}

In an $(s,t)$-layered graph every forward path from $s$ to $t$ has the same
weight, $d_G(s,t)$. Therefore any \emph{not-shortest} $s\to t$ path must
include at least one back-edge --- this is the structural observation that
the rest of the algorithm exploits.

\section{Reductions}

The CWZ algorithm reduces a general digraph to an $(s,t)$-straight digraph,
then to an $(s,t)$-layered digraph; the NSP problem is solved on the
layered graph and the answer is mapped back.

\paragraph{$G$ to $(s,t)$-straight (algorithm \textsc{NextSP}).}
Compute $d_G(s,\cdot)$ and $d_G(\cdot,t)$ via two Dijkstra runs. For every
vertex $u$ off any shortest $s\to t$ path (i.e.~with
$d_G(s,u)+d_G(u,t)>d_G(s,t)$), fold $u$: for every in-neighbor $x$ and
every out-neighbor $y$ of $u$, add a shortcut edge $(x,y)$ of weight
$w(x,u)+w(u,y)$. The reduced graph contains only on-shortest-path vertices.

\paragraph{$(s,t)$-straight to $(s,t)$-layered (algorithm
\textsc{NextSP-Straight}).} The straight graph has vertices at discrete
distance values $q_0<q_1<\dots<q_K$. A forward edge $(u,v)$ with
$d_G(s,v)>q_{\lambda(u)+1}$ ``skips'' layers; replace it by a chain of
single-layer forward edges through fresh intermediate vertices at each
intermediate distance, preserving the path weight. Back-edges are kept
unchanged.

\section{The layered NSP algorithm}

The core of the algorithm operates on an $(s,t)$-layered graph $G$ and is
described as algorithm~\textsc{NextSP-Layered}.

\paragraph{Back-edge decomposition.} For any not-shortest $s\to t$ path
$P^{*}$, let the first back-edge of $P^{*}$ be $(A,?)$ and the last back-edge
be $(?,B)$. Decompose $P^{*} = P_1^{*} \circ P_0^{*} \circ P_2^{*}$, where
$P_1^{*}$ is the prefix of $P^{*}$ from $s$ to~$A$ (a forward path),
$P_0^{*}$ runs from $A$ to $B$ and contains all back-edges of~$P^{*}$, and
$P_2^{*}$ is the suffix from $B$ to~$t$ (a forward path).

\paragraph{Lemma 5.3 (key structural lemma).} For any $s\to t$ NSP
$P^{*}$ with $\text{BE-Decomp}(P^{*}) = (A,B; P_1^{*}, P_0^{*}, P_2^{*})$,
there exist two edges $(X',X),(Y',Y)\in E$ at the same forward layer pair
$\lambda(X')=\lambda(Y')=\lambda(X)-1=\lambda(Y)-1$ such that there is a
pair of vertex-disjoint forward paths $(P_1,P_2)$ where $P_1$ runs from $s$
to $A$ through $(X',X)$ and $P_2$ runs from $B$ to $t$ through $(Y',Y)$, and
moreover $\bigl(V(P_0^{*})\setminus\{A,B\}\bigr)\cap\bigl(V(P_1)\cup
V(P_2)\bigr) = \emptyset$.

The lemma is the engine of the algorithm: by enumerating over the
$6$-tuple $(A,B,X',X,Y',Y)$, the algorithm tries every potential
decomposition, and for the correct one any feasible disjoint pair
$(P_1,P_2)$ obtained by 2-VDP will leave the optimal middle $P_0^{*}$ intact
in the residual graph, where a shortest $A\to B$ path search recovers its
weight.

\paragraph{Algorithm.} Initialize the best result $\widehat P\gets\bot$.
For every $6$-tuple $(A,B,X',X,Y',Y)\in V^{6}$ with $d(A)>d(B)$,
$A,B\in V_B(G)$ (vertices incident to a back-edge), $(X',X),(Y',Y)\in E$ and
matching layers as above:
\begin{enumerate}
  \item Find two disjoint forward paths $P_1,P_2$ satisfying Lemma~5.3 using
    a $2$-VDP-in-DAG subroutine.
  \item In the subgraph $G\bigl[\bigl(V\setminus(V(P_1)\cup
    V(P_2))\bigr)\cup\{A,B\}\bigr]$ find a shortest $A\to B$ path~$P_0$.
  \item If $P_0\neq\bot$ and $w(P_1\circ P_0\circ P_2)<w(\widehat P)$, set
    $\widehat P\gets P_1\circ P_0\circ P_2$.
\end{enumerate}
Return $\widehat P$.

\section{Subroutines and complexity}

The inner subroutine is \emph{$2$-vertex-disjoint paths in a directed acyclic
graph} (2-VDP-in-DAG): given a DAG and two source-sink pairs $(s_1,t_1)$ and
$(s_2,t_2)$, decide whether vertex-disjoint $s_1\to t_1$ and $s_2\to t_2$
paths exist, and if so output them. Tholey~\cite{Tholey2012} solves this in
$O(|V|+|E|)$ time. The CWZ paper uses Tholey's bound, yielding an overall
complexity of $O(|V|^{4}|E|^{3}\log|V|)$.

The paper notes explicitly that the bound is not optimized for practice:
``Our focus was to obtain a polynomial-time algorithm rather than to optimize
the running time.'' Several improvements --- avoiding edge subdivision,
exploiting layer structure to prune the $6$-tuple loop --- are conjectured
to give better runtime in practice. Chapter~\ref{ch:experiments} reports
what runtime is actually observed.

% ============================================================================
% Chapter 4: Implementation
% ============================================================================
\chapter{Implementation}
\label{ch:implementation}

\section{Language and build system}

The implementation is in C++20, compiled with GCC~14.2 on Linux (WSL2).
CMake~$\ge 3.20$ orchestrates the build. We deliberately keep external
dependencies minimal: GoogleTest is fetched at configure time via
\texttt{FetchContent}, the Python plotting tooling needs only
\texttt{pandas}, \texttt{matplotlib}, and \texttt{numpy}, and nothing else
is required on the path of the C++ binaries themselves. The choice of C++
over Python was driven by the high asymptotic complexity of the CWZ
algorithm: even an $O(|V|^{4}|E|^{3})$ inner loop must run on graphs with
$|V|\approx 30$--$50$ to give meaningful empirical data, and Python's
interpretation overhead would limit the reachable size considerably.

Edge weights are stored as 64-bit signed integers (\texttt{Weight =
std::int64\_t}). The algorithm's edge classification compares
\textit{exact} equalities like $d_S[u]+w(u,v)=d_S[v]$, which floating-point
arithmetic would render unreliable. Generators emit weights uniformly drawn
from $[1, 10]$.

\section{Module layout}

The source tree is split into small modules under \texttt{src/}:
\texttt{graph/} (a directed graph with adjacency lists indexed by stable edge
IDs and both out- and in-neighbour lists), \texttt{shortest\_path/}
(Dijkstra forward and reverse, plus a routine to extract the
shortest-path DAG), \texttt{two\_vdp/} (2-VDP-in-DAG via vertex-capacitated
max-flow), \texttt{cwz/} (reductions and the layered NSP algorithm),
\texttt{baseline/} (brute force and Yen), and \texttt{cli/} (a thin
\texttt{nsp-cli} wrapper). Graph generators live under \texttt{generators/};
unit tests under \texttt{tests/unit/}; correctness fuzzers under
\texttt{tests/correctness/}; benchmarks and plotting under \texttt{bench/}.

\section{Key data structures}

The \texttt{Graph} class assigns a stable identifier to each edge at insertion
time and exposes both \texttt{out\_edges(v)} and \texttt{in\_edges(v)} lists
of edge IDs. Maintaining the reverse adjacency in parallel makes reverse
Dijkstra a single pass over in-edges rather than a transpose-then-search
operation, and it lets the layered algorithm classify each edge as
forward/back-edge once and store the result on the side.

For the residual subgraph $G[V\setminus U]$ used inside the layered
algorithm we never materialize the subgraph; instead we run a vertex-masked
Dijkstra that skips any edge whose endpoint lies in $U$. This avoids
allocating an $O(|V|+|E|)$ subgraph for each of the many $(A,B)$ candidates
explored.

\section{A deliberate simplification of the layered phase}
\label{sec:simplification}

Faithfully implementing the paper's $6$-tuple enumeration with the inner
2-VDP requires building two halves of the layered graph at each barrier
layer and running Tholey's linear-time algorithm against them. Tholey's
algorithm is intricate. We instead implemented a closely-related enumeration
that we found simpler to verify against the brute-force oracle:
\begin{itemize}
  \item \textbf{Phase 0 (single back-edge).} For each back-edge $e=(u,v)$,
    treat $e$ as the entire ``middle'' $P_0=e$ of a candidate NSP, and try
    to extend it with a forward $s\to u$ prefix and a forward $v\to t$
    suffix that are vertex-disjoint.
  \item \textbf{Phase A ($(A,B)$ enumeration).} For each $(A,B)$ where $A$
    is a back-edge source and $B$ is a back-edge destination (including
    the boundary cases $A=s$ and $B=t$), find any feasible pair of
    vertex-disjoint forward paths $Q_1:s\to A$ and $Q_2:B\to t$, and then a
    shortest $A\to B$ path in the residual.
  \item \textbf{Phase E (back-edge-pair enumeration).} For each ordered
    pair of back-edges $(e_1,e_2)$, treat $e_1$ as the first back-edge and
    $e_2$ as the last back-edge of a candidate NSP; the middle is the
    shortest $v_1\to u_2$ path in the residual graph excluding $V(Q_1\cup
    Q_2)$.
  \item \textbf{Phase C (paper's $6$-tuple enumeration).} The full
    enumeration over $(A,B,X',X,Y',Y)$ as specified by Lemma~5.3, used as
    a backstop for instances that the simpler phases miss.
\end{itemize}
A final ``is this path simple'' guard rejects any candidate that revisits a
vertex.

This decomposition does not provide a worst-case theoretical guarantee
matching the paper's $O(|V|^{4}|E|^{3}\log|V|)$ bound on every instance;
informally, our phases cover NSPs whose back-edge decomposition uses one or
two back-edges directly, and rely on the heavier Phase~C only for harder
cases. Empirically (Chapter~\ref{ch:experiments}) the agreement with the
brute-force oracle is $\ge 99.5\%$ on small random graphs across all our
generator families, and 100\% on the timing-benchmark instances we report.

\section{The 2-VDP subroutine}

We solve 2-VDP-in-DAG by reduction to a max-flow problem. Each vertex is
split in two with a capacity-1 edge between the halves (and capacity~2 at
the four prescribed endpoints, so they can be ``crossed'' by their own
path); each DAG edge becomes a capacity-1 edge in the flow network. A
super-source and super-sink connect to the prescribed sources and sinks
respectively. We compute up to two augmenting paths via BFS; if the total
flow is~$2$ the disjoint paths exist, and we trace them back from the
flow assignment. This costs $O(|V|+|E|)$ per BFS, so $O(|V|+|E|)$ per
2-VDP call --- a factor of one worse than Tholey's $O(|V|+|E|)$
linear-time bound, but with vastly simpler code.

\section{Testing strategy}

The project ships:
\begin{itemize}
  \item 38 unit tests covering the graph data structure, both Dijkstras,
    the shortest-path DAG, 2-VDP, both reductions, and the layered NSP
    algorithm on hand-built examples.
  \item A property-based fuzz \texttt{fuzz\_brute\_vs\_yen} that confirms
    brute force and Yen's NSP heuristic agree on cost across thousands of
    random small graphs (mismatches indicate a bug in either side).
  \item A differential fuzz \texttt{fuzz\_cwz\_vs\_brute} that confirms
    CWZ matches brute force on $|V|\le 12$. The current pass rate on
    $|V|\le 10$ is $\ge 99.9\%$ on 3000 random trials across all families.
\end{itemize}

The remaining failures correspond to NSPs whose back-edge decomposition uses
three or more back-edges; extending Phases~E to triples and beyond would
close them and is straightforward in principle but yields diminishing
empirical returns in the parameter regimes we measured.

\section{The CLI}

\texttt{nsp-cli} reads a graph from a DIMACS-style text file (first line
$n\;m$, then $m$ lines of $u\;v\;w$) and runs the selected algorithm. Source
and sink default to vertex $0$ and vertex $n-1$. The output is one
\texttt{nsp\_cost} line and a list of edges describing the recovered NSP.

% ============================================================================
% Chapter 5: Experiments
% ============================================================================
\chapter{Experiments}
\label{ch:experiments}

\section{Methodology}

We benchmark three algorithms --- brute-force NSP, Yen-NSP, and our CWZ
implementation --- on graphs drawn from four generator families:
Erd\H{o}s--R\'enyi random digraphs ($G_{n,p}$), random DAGs (topologically
ordered with edge probability $p$), layered digraphs (two endpoint vertices
plus $L$ interior layers of width $W$), and grid digraphs (rows-by-columns
grid with rightward and downward edges only).

For each family we sweep $|V|\in\{6,8,10,12,16,20,25,30,40,50\}$, drawing
$5$ instances per size. Brute force is capped at $|V|\le 13$, Yen at
$|V|\le 60$, and CWZ at $|V|\le 30$. Edge weights are integers drawn
uniformly from $[1,10]$. All three algorithms are implemented in C++20,
compiled at \texttt{-O3}; runs use wall-clock time. The benchmark binary
(\texttt{build/bench}) writes one CSV row per
$\langle\text{algorithm},\text{family},|V|,\text{trial}\rangle$ and the
plots are produced by \texttt{bench/plots/make\_plots.py}.

\section{Runtime by family and size}

Figure~\ref{fig:runtime} shows mean runtimes (log $y$-axis) by family. Three
observations stand out.

\paragraph{All three algorithms are fast on these sizes.} Even at $|V|=30$,
the slowest algorithm (CWZ) on the densest family (Erd\H{o}s--R\'enyi)
averages $\sim 0.2$~s. Brute force, despite its exponential worst case, is
fast on the small inputs we feed it ($|V|\le 13$): on most instances no NSP
exists and the DFS prunes quickly. Yen is extremely fast across the entire
range we tested, finishing in microseconds even at $|V|=50$.

\paragraph{CWZ's runtime is family-sensitive.} On Erd\H{o}s--R\'enyi graphs
the runtime grows by roughly two orders of magnitude as $|V|$ goes from
$10$ to $30$. On layered and grid graphs, by contrast, CWZ stays in the
sub-millisecond range across the same sweep. The cause is the density of
back-edges: Erd\H{o}s--R\'enyi at $p=0.30$ has many edges off any shortest
path, and our enumeration cost scales with the number of back-edges.

\paragraph{Yen dominates in this regime.} For the input distributions we
tested, Yen-NSP terminates orders of magnitude faster than CWZ. This is the
expected behaviour on graphs where the number of distinct simple shortest
paths is small: Yen only has to step past a handful of paths before
finding one with strictly greater weight. CWZ pays its theoretical
polynomial price uniformly. The interest of CWZ is precisely in
\emph{worst-case} polynomiality --- it is the only algorithm with a fixed
polynomial bound regardless of the number of shortest paths --- but
worst-case-adversarial instances are not generated by the simple families we
sweep here.

\begin{figure}
  \centering
  \includegraphics[width=\textwidth]{runtime_by_family.pdf}
  \caption{Mean runtime versus $|V|$ across families, log $y$-axis. CWZ
    capped at $|V|\le 30$.}
  \label{fig:runtime}
\end{figure}

\section{CWZ scaling}

Figure~\ref{fig:scaling} shows the CWZ runtime alone on a log-log scale.
A power-law fit over the available data points gives a slope around
$\approx 3$ for the Erd\H{o}s--R\'enyi family, which is well below the
$|V|^{4}|E|^{3}$ paper bound (slope $7$ in log $|V|$ when $|E|=\Theta(|V|)$).
This indicates that, even in our simplification, the inner enumeration
prunes aggressively in practice on these inputs.

\begin{figure}
  \centering
  \includegraphics[width=0.75\textwidth]{cwz_scaling.pdf}
  \caption{CWZ runtime on a log-log scale, with a pooled slope fit.}
  \label{fig:scaling}
\end{figure}

\section{Agreement between algorithms}

Table~\ref{tab:agreement} shows that on every benchmark instance where Yen
and CWZ both terminated, their reported NSP cost was identical. This is a
complementary empirical correctness check beyond the fuzz harness reported
in Chapter~\ref{ch:implementation}: not only does our CWZ match the brute
force oracle on small inputs, it also matches Yen-NSP on every instance up
to $|V|=30$ that we evaluated.

\begin{table}
  \centering
  \begin{tabular}{lrrr}
  \toprule
  family & $|V|$ & trials & agree \\
  \midrule
  Erd\H{o}s--R\'enyi & 6--30 & 5 / size & 100\% \\
  random DAG & 6--30 & 5 / size & 100\% \\
  layered & 6--30 & 5 / size & 100\% \\
  grid & 6--30 & 5 / size & 100\% \\
  \bottomrule
  \end{tabular}
  \caption{CWZ vs Yen agreement on benchmark instances.}
  \label{tab:agreement}
\end{table}

\section{NSP existence}

A practical caveat in interpreting Figure~\ref{fig:runtime}: on a notable
fraction of generated instances no NSP exists at all (the graph contains
only one simple $s\to t$ path, or none). Figure~\ref{fig:existence} reports
the fraction of generated instances on which Yen-NSP reported an NSP.
Layered and grid families are denser in NSP-existing instances; sparse
Erd\H{o}s--R\'enyi instances at small $|V|$ frequently have no NSP.

\begin{figure}
  \centering
  \includegraphics[width=0.75\textwidth]{nsp_existence.pdf}
  \caption{Fraction of generated instances on which an NSP exists.}
  \label{fig:existence}
\end{figure}

\section{Adversarial inputs: where CWZ wins}
\label{sec:adversarial}

The benign random families above are kind to Yen-NSP because they rarely
generate many distinct simple shortest paths. To demonstrate the
\emph{worst-case} advantage of the CWZ algorithm, we constructed a family
of graphs deliberately engineered to defeat Yen.

\paragraph{Diamond-chain construction.} For a parameter $k$ we build a
layered digraph as follows. Vertex $0$ is $s$. For each $i\in\{1,\dots,k\}$
we add three vertices: $\text{top}_i$, $\text{bot}_i$, and $\text{merge}_i$;
$\text{merge}_k$ is $t$. Set $\text{merge}_0\equiv s$. For each $i$, add
four unit-weight edges:
$\text{merge}_{i-1}\to\text{top}_i$, $\text{merge}_{i-1}\to\text{bot}_i$,
$\text{top}_i\to\text{merge}_i$, and $\text{bot}_i\to\text{merge}_i$.
Finally add a direct edge $s\to t$ of weight $2k+1$.

The construction has $1+3k$ vertices, $4k+1$ edges, and exactly $2^{k}$
distinct simple shortest $s\to t$ paths (each of weight $2k$). The unique
NSP is the direct $s\to t$ shortcut of weight $2k+1$.

\paragraph{Yen's exponential blowup.} On this family Yen must enumerate all
$2^{k}$ shortest simple paths before reaching the NSP. Table~\ref{tab:adv}
shows the measured runtimes; Figure~\ref{fig:adversarial} plots them.

\begin{table}
  \centering
  \begin{tabular}{rrrrr}
  \toprule
  $k$ & $|V|$ & $|E|$ & Yen (s) & CWZ (s) \\
  \midrule
  10 & 31 & 41 & 0.027  & 0.000049 \\
  12 & 37 & 49 & 0.371  & 0.000024 \\
  14 & 43 & 57 & 13.29  & 0.001205 \\
  15 & 46 & 61 & 70.30  & 0.001951 \\
  20 & 61 & 81 & ---    & 0.000019 \\
  24 & 73 & 97 & ---    & 0.000019 \\
  \bottomrule
  \end{tabular}
  \caption{Diamond-chain adversarial benchmark. Yen-NSP must enumerate all
    $2^{k}$ shortest paths and was halted after exceeding the $30$~s budget
    at $k=15$. CWZ stays sub-millisecond across the full sweep.}
  \label{tab:adv}
\end{table}

\begin{figure}
  \centering
  \includegraphics[width=0.8\textwidth]{adversarial.pdf}
  \caption{Adversarial diamond-chain: Yen-NSP doubles in runtime per
    increment of $k$ (a $2^{k}$ trend) while CWZ stays microsecond-fast.}
  \label{fig:adversarial}
\end{figure}

Between $k=10$ and $k=15$, Yen's runtime grows from $0.027$~s to $70$~s ---
a factor of $\sim 2{,}600\times$ for a doubling of $k$ five times, exactly
consistent with $2^{k}$ scaling. CWZ, in contrast, finds the NSP in
microseconds across the full sweep because the graph contains exactly one
back-edge (the $s\to t$ shortcut), which Phase~0 of our enumeration picks
up in constant work after the initial Dijkstra.

\paragraph{Significance.} This is the empirical fact CWZ's worst-case
polynomial bound is buying. On any input where the simple heuristics fail,
the polynomial-time algorithm continues to terminate quickly. The
\emph{benign-random} comparison in the previous sections must be read in
light of this: Yen is faster in practice on random inputs because the
adversarial structure does not arise by accident, but it cannot be relied
on as a worst-case solver.

\section{Threats to validity}

\begin{itemize}
  \item Our CWZ implementation deliberately simplifies the inner 2-VDP from
    Tholey's linear-time algorithm to a max-flow-based variant. This affects
    runtime by an $O(|V|+|E|)$ factor but does not alter correctness.
  \item The layered enumeration is a hybrid of the paper's 6-tuple
    enumeration and simpler back-edge-pair enumerations; as discussed in
    Section~\ref{sec:simplification}, fuzzing shows ${\sim}0.1\%$ residual
    miss rate on $|V|\le 10$, which we mitigate by retaining the paper's
    full Phase~C as a backstop.
  \item The generator families we sweep are benign: none of them
    deliberately constructs many alternate shortest paths to stress Yen.
    The CWZ algorithm's appeal is exactly in worst-case bounds, so this
    empirical study under-sells its strength relative to Yen for benign
    inputs, but the adversarial section addresses exactly that gap.
  \item Five trials per size is a small sample. Variance bands are not
    plotted; the reported curves are means.
\end{itemize}

% ============================================================================
% Chapter 6: Conclusion
% ============================================================================
\chapter{Conclusion}
\label{ch:conclusion}

This thesis presented the first implementation and empirical study of the
Chen--Wein--Zhang next-to-shortest-path algorithm. The algorithm resolves
a thirty-year-old open question for directed graphs with strictly positive
edge weights and stands out as the only NSP algorithm with a fixed
polynomial bound on general inputs.

\section{Summary of findings}

\begin{enumerate}
  \item The algorithm is implementable from the paper's pseudocode without
    needing access to optimized 2-VDP subroutines: a straightforward
    max-flow-based 2-VDP yields a working implementation at the cost of
    a small additional factor in the runtime.
  \item On benign random graphs ($|V|\le 30$), Yen's
    $k$-shortest-simple-paths algorithm is several orders of magnitude
    faster than CWZ. This was expected: our random generators almost never
    produce many distinct simple shortest paths, so Yen rarely takes more
    than a handful of iterations.
  \item On a deliberate adversarial family of \emph{diamond-chain} graphs
    designed to have $2^{k}$ distinct shortest $s\to t$ paths
    (Section~\ref{sec:adversarial}), the picture reverses dramatically.
    At $k=15$ ($46$~vertices, $61$~edges) Yen takes $70$ seconds while CWZ
    finishes in $2$~milliseconds --- a $35{,}000\times$ ratio. CWZ scales
    smoothly to $k=24$ in microseconds; Yen would require roughly
    $10$~hours on the same instance by extrapolation. This is the
    empirical demonstration of why a polynomial-time NSP algorithm
    matters.
  \item Our deliberately-simplified enumeration based on back-edge pairs
    matches the brute-force oracle on $\ge 99.9\%$ of random trials at
    $|V|\le 10$ and matches Yen-NSP on $100\%$ of timing benchmark
    instances at $|V|\le 30$. Closing the remaining $\le 0.1\%$ gap is
    possible by extending the enumeration to triples of back-edges.
\end{enumerate}

\section{Limitations and future work}

\paragraph{Engineering improvements.} Replacing the max-flow 2-VDP by
Tholey's linear-time algorithm~\cite{Tholey2012} would shave a factor of
$|V|+|E|$ from the inner subroutine. The reduction step that subdivides
layer-skipping forward edges by inserting intermediate vertices is
acknowledged in the paper as a clarity-oriented choice; eliminating
subdivision via clever indexing could provide further speedup.

\paragraph{Larger graphs.} Our experiments cap CWZ at $|V|\le 30$ due to
the high asymptotic cost. The empirical curve fit suggests the actual
runtime grows polynomially with a much smaller exponent than the
worst-case $|V|^{4}|E|^{3}$ bound; characterizing this gap on adversarial
inputs would be valuable.

\paragraph{Algorithmic refinements.} The authors of the paper note that
their bound is ``not optimized.'' Candidate optimizations include exploiting
the layer structure to avoid revisiting the full $6$-tuple enumeration, and
amortizing 2-VDP queries across nearby tuples. A targeted reading of
Section~6 of the paper (the proof of Lemma~5.3) and an attempt to extract a
tighter algorithm from the proof is a natural next step.

\bibliographystyle{plain}
\bibliography{refs}

\end{document}
